\section*{Capítulo \ref{ch:g11:plate-motions} --- As placas e os seus movimentos}

\begin{solution}{exo:g11:plate-motions:1}
Uma \omterm{def:g11:plate-motions:rigid-plate}{placa rígida} é um pedaço de \omterm{def:g11:lithosphere-asthenosphere:lithosphere}{litosfera} que se move como um único bloco:
as distâncias entre os seus pontos não mudam, e a sua deformação se
concentra em zonas estreitas ao longo dos seus limites. Os terremotos
mostram isso: quase todos acontecem ao longo de faixas estreitas, os limites
de placas, e muito poucos no interior das placas.
\end{solution}

\begin{solution}{exo:g11:plate-motions:2}
\qty{72}{mm} por ano equivalem a \qty{72}{km} por milhão de anos. Em 30
anos, a estação se desloca $72 \times 30 = \qty{2160}{mm}$, cerca de
\qty{2.2}{m}; em um milhão de anos, \qty{72}{km}.
\end{solution}

\begin{solution}{exo:g11:plate-motions:3}
$v_E = 30/10 = \qty{3}{cm}$ por ano e $v_N = 40/10 = \qty{4}{cm}$ por ano. A
velocidade é $\sqrt{3^2 + 4^2} = \qty{5}{cm}$ por ano. O ângulo $\beta$ com
o norte satisfaz $\tan\beta = 3/4$, logo $\beta \approx \ang{37}$: a estação
se move para nordeste, \ang{37} a leste do norte.
\end{solution}

\begin{solution}{exo:g11:plate-motions:4}
Um \omterm{def:g11:plate-motions:hotspot-track}{rastro de ponto quente} é a linha de vulcões e de montes submarinos que
uma placa leva para longe de um ponto quente ao se mover sobre ele, cada vez
mais antigos com a distância ao vulcão ativo. O vulcão mais jovem, o
Kilauea, fica na extremidade sudeste da cadeia, na ilha do Havaí; as idades
aumentam para oés-noroeste. A placa leva os seus vulcões para longe do ponto
quente no sentido do seu movimento: a placa do Pacífico se move para
oés-noroeste.
\end{solution}

\begin{solution}{exo:g11:plate-motions:5}
$519/5.1 \approx \qty{102}{km}$ por milhão de anos, isto é, \qty{102}{mm}
por ano, cerca de \qty{10}{cm} por ano.
\end{solution}

\begin{solution}{exo:g11:plate-motions:6}
$v_E = 63/31 \approx \qty{20.3}{mm}$ por ano e $v_N = 48/31 \approx
\qty{15.5}{mm}$ por ano. A velocidade é $\sqrt{20.3^2 + 15.5^2} \approx
\qty{25.5}{mm}$ por ano; $\tan\alpha = 15.5/20.3 \approx 0.76$ dá $\alpha
\approx \ang{37}$: a estação se move \ang{37} ao norte do leste, para
nordeste. O Mauna Kea se move $72/25.5 \approx 2.8$ vezes mais depressa.
\end{solution}

\begin{solution}{exo:g11:plate-motions:7}
(1) Em milímetros por ano,
\[
  \vec v_{\mathrm{NA/EU}} = (-10.5 - 7.7,\ 20.6 - 16.3) = (-18.2, +4.3).
\] A sua velocidade escalar é $\sqrt{18.2^2 + 4.3^2} \approx
\qty{18.7}{mm}$ por ano, e $\tan\alpha = 4.3/18.2 \approx 0.24$ dá
\ang{13} ao norte do oeste. (2) As duas placas se afastam através da dorsal
a cerca de \qty{1.9}{cm} por ano: é a \omterm{def:g11:drift-to-spreading:spreading-rate}{taxa de expansão total} da dorsal na
Islândia. (3) Isso concorda com os \qty{2}{cm} por ano citados para o
Atlântico Norte.
\end{solution}

\begin{solution}{exo:g11:plate-motions:8}
(1) $v = \pi \times 0.651 \times 6371/180 \approx \qty{72.4}{km}$ por milhão
de anos, cerca de \qty{7.2}{cm} por ano. (2) $72.4 \times \sin\ang{30} \approx
\qty{36.2}{km}$ por milhão de anos, cerca de \qty{3.6}{cm} por ano. (3) No
polo de rotação, se ele fica sobre a placa. A velocidade muda de ponto para
ponto, de zero no polo até o seu valor máximo a \ang{90} dele: uma placa
inteira não tem uma única velocidade, apenas uma velocidade em cada ponto
(ou uma média sobre a sua superfície).
\end{solution}

\begin{solution}{exo:g11:plate-motions:9}
(1) Em milímetros por ano,
\[
  \vec v_{\mathrm{IN/EU}} = (40.6 - 24.0,\ 38.6 + 3.3) = (+16.6, +41.9).
\] A velocidade é $\sqrt{16.6^2 + 41.9^2} \approx
\qty{45}{mm}$ por ano; o ângulo $\beta$ com o norte satisfaz $\tan\beta =
16.6/41.9 \approx 0.40$, logo $\beta \approx \ang{22}$: a Índia se move em
direção à Eurásia, \ang{22} a leste do norte. (2) \qty{45}{km} por milhão de
anos, logo \qty{450}{km} em 10 milhões de anos.
\end{solution}

\begin{solution}{exo:g11:plate-motions:10}
$(2432 - 1058)/(27.7 - 10.3) = 1374/17.4 \approx \qty{79}{km}$ por milhão
de anos, cerca de \qty{7.9}{cm} por ano, um pouco menos que \qty{9.1}{cm}
por ano. As idades têm cada uma uma incerteza de cerca de meio milhão de
anos; a diferença de idade, $17.4 \pm 1.0$ milhões de anos, dá uma
velocidade entre cerca de 75 e \qty{84}{km} por milhão de anos. Entre esses
dois vulcões, a placa parece ter se movido um pouco mais devagar que em
média.
\end{solution}

\begin{solution}{exo:g11:plate-motions:11}
(1) $\vec v_{A/C} = \vec v_{A/B} + \vec v_{B/C}$: \qty{3}{cm} por ano para
leste mais \qty{4}{cm} por ano para o norte. As duas setas formam um ângulo
reto, e o terceiro lado do triângulo tem o comprimento
$\sqrt{3^2 + 4^2} = \qty{5}{cm}$ por ano, na direção \ang{37} a leste do
norte ($\tan\beta = 3/4$). (2) $\vec v_{C/A}
= -\vec v_{A/C}$: \qty{5}{cm} por ano no sentido oposto, \ang{37} a oeste
do sul.
\end{solution}

\begin{solution}{exo:g11:plate-motions:12}
(1) $54/0.375 = \qty{144}{km}$ por milhão de anos, \qty{14.4}{cm} por ano.
(2) De $29/0.375 \approx 77$ a $79/0.375 \approx \qty{211}{km}$ por milhão
de anos: entre cerca de 8 e \qty{21}{cm} por ano. (3) Perto do ponto
quente, a distância não é maior que o tamanho de um vulcão ou da região
sobre a qual o ponto quente alimenta vulcões, e a idade não é maior que o
tempo que um vulcão leva para crescer: as duas são muito incertas em termos
relativos. Longe do ponto quente, um erro de \qty{25}{km} em
\qty{2400}{km}, ou de meio milhão de anos em 28 milhões, é de apenas um ou
dois por cento.
\end{solution}

\begin{solution}{exo:g11:plate-motions:13}
(1) Dez vezes a dispersão dá \qty{50}{mm}: $50/72 \approx 0.7$ ano, cerca de
oito meses. (2) $50/5 = 10$ anos. (3) Ao longo de milhões de anos, as placas
se deslocam dezenas a centenas de quilômetros, distâncias que as listras
magnéticas registram e que eram fáceis de medir no mar; um posicionamento
com precisão de poucos milímetros, que torna um movimento visível em meses,
veio muito depois.
\end{solution}

\begin{solution}{exo:g11:plate-motions:14}
(1) Em milímetros por ano,
\[
  \vec v_{\mathrm{SA/AF}} = (-4.9 - 27.1,\ 11.0 - 17.9) = (-32.0, -6.9).
\] A velocidade é $\sqrt{32.0^2 + 6.9^2} \approx
\qty{32.7}{mm}$ por ano, e $\tan\alpha = 6.9/32.0 \approx 0.22$ dá
\ang{12} ao sul do oeste. (2) A \omterm{def:g11:drift-to-spreading:spreading-rate}{taxa de expansão total} é de cerca de
\qty{3.3}{cm} por ano, a semitaxa de cerca de \qty{1.6}{cm} por ano. (3) A
semitaxa dos últimos milhões de anos, \qty{1.6}{cm} por ano, é um pouco
menor que os \qty{1.9}{cm} por ano em média ao longo de 66 milhões de anos:
o Atlântico Sul pode ter se aberto um pouco mais devagar recentemente que em
média, ou a diferença pode vir das incertezas das idades dos sedimentos. As
duas medições concordam em que a dorsal é lenta.
\end{solution}

\begin{solution}{exo:g11:plate-motions:15}
(1) Uma reta que sobe \qty{20}{mm} por ano, uma queda brusca de
\qty{50}{mm} depois de cinco anos, e então uma reta que volta a subir com a
mesma inclinação. (2) A estação termina $10 \times 20 - 50 = \qty{150}{mm}$
a leste do ponto de partida: $150/10 = \qty{15}{mm}$ por ano, um quarto a
menos. (3) Cada metade dá \qty{20}{mm} por ano. (4) O salto é um movimento
separado e brusco, causado pelo terremoto: o geodesista mede o seu tamanho
pelas posições logo antes e logo depois dele, retira-o da série (ou ajusta
as duas partes com a mesma inclinação e um degrau entre elas), e mantém
\qty{20}{mm} por ano como a velocidade da estação.
\end{solution}

\begin{solution}{pb:g11:plate-motions:1}
\textbf{1.} O Mauna Kea (\num{0.375} milhão de anos) é o mais jovem, Suiko
(centro) (\num{64.7} milhões de anos) o mais antigo. As idades aumentam com a
distância ao Kilauea.

\textbf{2.} $\qty{1}{km} = 10^6~\unit{mm}$ e um milhão de anos são $10^6$
anos: \qty{1}{km} por milhão de anos é $10^6~\unit{mm}/10^6$ anos, isto é,
\qty{1}{mm} por ano.

\textbf{3.} $519/5.1 \approx \qty{102}{mm}$ por ano.

\textbf{4.} Necker $1058/10.3 \approx 102.7$; Laysan $1818/19.9 \approx
91.4$; Midway $2432/27.7 \approx 87.8$; Colahan $3128/38.6 \approx 81.0$;
Abbott $3280/38.7 \approx 84.8$ (milímetros por ano).

\textbf{5.} Os pontos (54; 0.375), (519; 5.1), (780; 7.2), (1058; 10.3),
(1818; 19.9), (2432; 27.7), (3128; 38.6), (3280; 38.7) ficam perto de uma
reta que passa pela origem, com alguma dispersão (em relação a uma reta de
\qty{91}{km} por milhão de anos, Kauai, Nihoa e Necker ficam um pouco
abaixo dela, Laysan sobre ela, Colahan e Abbott um pouco acima).

\textbf{6.} Cada vulcão foi construído sobre o ponto quente e desde então se
move com a placa. No tempo $t$, a placa se deslocou $v\,t$ numa direção
constante, logo o vulcão fica a $d = v\,t$ do ponto que agora está sobre o
ponto quente: $t = d/v$.

\textbf{7.} $(102.7 + 91.4 + 87.8 + 81.0 + 84.8)/5 \approx \qty{89.5}{mm}$ por
ano.

\textbf{8.} 89.5 contra \qty{91}{km} por milhão de anos: a diferença é
menor que \qty{2}{\percent}.

\textbf{9.} $519/91 \approx 5.7$ milhões de anos, contra uma idade de
\num{5.1} milhões de anos: perto, sendo Kauai um pouco mais jovem do que a
reta média prevê.

\textbf{10.} Os vulcões são cada vez mais antigos para oés-noroeste: a placa
os leva nessa direção. A placa do Pacífico se move para oés-noroeste sobre o
ponto quente.

\textbf{11.} $(4860 - 3493)/(64.7 - 42.4) = 1367/22.3 \approx \qty{61}{km}$
por milhão de anos, cerca de \qty{6.1}{cm} por ano.

\textbf{12.} $65/91 \approx 0.71$: ao longo da cadeia do Imperador, a placa
se moveu cerca de \qty{30}{\percent} mais devagar que ao longo da cadeia
havaiana.

\textbf{13.} $3668/50 \approx \qty{73}{km}$ por milhão de anos, cerca de
\qty{7.3}{cm} por ano, menos que \qty{91}{km} por milhão de anos. A reta
pela origem usa as idades de 1987, que situam o fim da cadeia havaiana entre
39 e 43 milhões de anos; a idade mais recente, de cerca de 50 milhões de
anos, para o início da curva espalha a mesma distância por um tempo mais
longo. Os dois estão certos se a placa se moveu mais devagar no início da
cadeia havaiana e mais depressa depois: a velocidade não foi constante.

\textbf{14.} Uma mudança na direção do movimento da placa do Pacífico, há
cerca de 50 milhões de anos.

\textbf{15.} A direção da cadeia é a direção em que a placa se moveu; o
espaçamento das idades ao longo dela é a sua velocidade. Uma placa pode
acelerar ou desacelerar sem virar, ou virar sem mudar de velocidade, de modo
que as duas informações são independentes. Mudanças de velocidade sem
mudança de direção ao longo da cadeia do Imperador mostram que o seu
registro de velocidade e a curva são eventos separados: a curva não foi
produzida por uma desaceleração em relação ao ponto quente.

\textbf{16.} $\sqrt{63.2^2 + 34.5^2} \approx \qty{72.0}{mm}$ por ano, e
$\tan\alpha = 34.5/63.2 \approx 0.55$ dá \ang{29} ao norte do oeste.

\textbf{17.} Oés-noroeste fica a \ang{22.5} ao norte do oeste, noroeste a
\ang{45}: a direção da estação, \ang{29} ao norte do oeste, fica entre as
duas, perto da direção da cadeia. Ela é compatível.

\textbf{18.} A velocidade da estação é medida em relação a um referencial em
que a \omterm{def:g11:lithosphere-asthenosphere:lithosphere}{litosfera} como um todo não gira, enquanto a cadeia dá o movimento em
relação ao ponto quente: se os dois referenciais se movem um em relação ao
outro, as velocidades diferem. E a estação mede os últimos 30 anos, a cadeia
os últimos 40 milhões de anos, durante os quais a velocidade pode ter
mudado. A dispersão das idades também deixa uma incerteza de alguns
milímetros por ano sobre a velocidade da cadeia.

\textbf{19.} $100/91 \approx 1.1$ milhão de anos.

\textbf{20.} \textbf{A placa do Pacífico se moveu sobre o ponto quente do
Havaí a cerca de \qty{9}{cm} por ano} durante os últimos 40 milhões de anos,
para oés-noroeste; antes da curva, ela se movia mais devagar, e em outra
direção.
\end{solution}
